\subsection{Adding actual harmonic functions}\label{subsec:fc-line-subharmonic-harmonic-addition}

\begin{definition}[Harmonicity on an open line domain]\label{def:fc-actual-line-harmonic-on}
\lean{EnvelopeProperty.Berkovich.RelativeSpectrum.IsLineHarmonicOn}
\leanok
\uses{def:fc-actual-harmonic}
A real function $h$ is harmonic on $\Omega$ if its restriction to every
actual standard affinoid contained in $\Omega$ is an actual harmonic function
in the logarithmic-coordinate sense.
\end{definition}

\begin{lemma}[Continuity of actual harmonic functions on open domains]\label{lem:fc-line-harmonic-on-continuity}
\lean{EnvelopeProperty.Berkovich.RelativeSpectrum.IsLineHarmonicOn.continuousOn, EnvelopeProperty.Berkovich.RelativeSpectrum.IsLineHarmonicOn.const, EnvelopeProperty.Berkovich.RelativeSpectrum.IsLineHarmonicOn.mono}
\leanok
\uses{def:fc-actual-line-harmonic-on, lem:fc-harmonic-retraction,
  thm:fc-standard-affinoid-basis}
On an open domain $\Omega$, an actual harmonic function is continuous.
Constants are harmonic, and harmonic functions restrict to smaller domains.
\end{lemma}

\begin{proof}
Every point of $\Omega$ belongs to the interior of an actual standard affinoid
contained in $\Omega$. The given function is actually harmonic on that
affinoid, hence continuous there. A set containing a neighbourhood of the
point therefore witnesses continuity at the point. This proves continuity
on $\Omega$. Constants and restriction follow directly from the literal
finite-coordinate definition on each test affinoid.
\end{proof}

\begin{lemma}[Addition of actual harmonic functions]\label{lem:fc-line-subharmonic-harmonic-addition}
\lean{EnvelopeProperty.Berkovich.RelativeSpectrum.IsLineSubharmonicOn.add_harmonic, EnvelopeProperty.Berkovich.RelativeSpectrum.IsLineSubharmonicOn.add_harmonic_of_continuousOn, EnvelopeProperty.Berkovich.RelativeSpectrum.IsLineSubharmonicOn.add_const}
\leanok
\uses{def:fc-actual-line-subharmonic, def:fc-actual-line-harmonic-on,
  lem:fc-line-harmonic-on-continuity, lem:fc-harmonic-vector-space}
If $u$ is subharmonic and $h$ is harmonic on the open domain $\Omega$, then
$u+h$ is subharmonic. In particular, addition of a real constant preserves
subharmonicity.
\end{lemma}

\begin{proof}
Continuity of the finite real function $h$ and upper semicontinuity of $u$
make $u+h$ upper semicontinuous; addition of a finite real number preserves
exclusion of $+\infty$. On any actual test affinoid $U$, if $u+h\le k$ on
the boundary for actual harmonic $k$, then $u\le k-h$ on that boundary.
The actual function $k-h$ is harmonic by vector-space closure. Comparison
for $u$ therefore gives $u\le k-h$ throughout $U$, which is equivalent to
$u+h\le k$. For a constant $h$, its continuity and harmonicity hold on any
domain, even without the openness condition.
\end{proof}
